% Pursuing Stacks (cotes 134-2 to 134-8): archive cote 134-6, site page 6 (PDF page 7), typescript page 325.
% typescript. Draft compared: ps5.tex:1621-1699.
% AI-transcribed by Claude Opus 5.5 (claude-opus-5-5, adaptive thinking, effort high), Message Batches API, batch ps3 (msgbatch_01XdFmuVVBAcgc7EFxdtCqsj), 2026-11-12; reviewed against both supplied scan halves; see CHECK_DECISIONS.tsv.
\scanpage{134-6/6}
\mspage{325}

(11) \qquad $j^{*} : \overset{\deleted{\unsure{\wedge}}}{\over{A}{A}}^{\wedge}\hand{\longrightarrow} \over{B}{B}^{\wedge}$

the\over{n}{n} satisfies
\[ j^{*\,-1}(\underline{W}_B) = \underline{W}_A \quad , \]
and the corresponding functor for the localizations gives rise to a
commutative diagram
% Hand-drawn arrows are marked blue; typed objects remain plain.
\diagram{\begin{tikzcd}[column sep=small]
\mathrm{Hot}_A \ar[color=blue,rr,"\hand{\overline{j^*}}"] \ar[color=blue,dr,"\hand{\overline{i}_A}"'] \& \& \mathrm{Hot}_B \ar[color=blue,dl,"\hand{\overline{i}_B}"] \\
\& (\mathrm{Hot}) \&
\end{tikzcd} \quad ,}
and hence the corresponding diagram

(12)
% Hand-drawn arrows are marked blue; typed objects remain plain.
\diagram{\begin{tikzcd}[column sep=small]
\mathrm{Hot}_A \ar[color=blue,rr,"\hand{\overline{j^*}}"] \ar[color=blue,dr] \& \& \mathrm{Hot}_B \ar[color=blue,dl] \\
\& \mathrm{D.(Ab)=Hotab} \&
\end{tikzcd}}
is commutative, where the vertical arrows are the canonical functors
(6). Coming back to the case B=\hand{$\Delta$} say, this shows that (6) can be
viewed as the composition
\[ \mathrm{Hot}_A \overset{\hand{\overline{j^*}}}{\hand{\longrightarrow}} \mathrm{Hot}_{\hand{\Delta}} \overset{\hand{\overline{Wh}_{\Delta}}}{\hand{\longrightarrow}} \mathrm{Hotab}_{\hand{\Delta}} \overset{\substack{\hand{\overline{DP}}\\ \hand{\approx}}}{\hand{\longrightarrow}} \mathrm{Hotab} \quad , \]
and hence it can be inserted in the commutative diagram

(13)
% Hand-drawn arrows are marked blue; typed objects remain plain.
\diagram{\begin{tikzcd}
\mathrm{Hot}_A \ar[color=blue,r] \ar[color=blue,d] \& \mathrm{Hot}_{\hand{\Delta}} \ar[color=blue,d] \& \\
\mathrm{Hotab}_A \ar[color=blue,r] \& \mathrm{Hotab}_{\hand{\Delta}} \ar[color=blue,r,"\hand{\approx}"] \& \mathrm{Hotab}
\end{tikzcd} \quad ,}
where the functor
\[ \mathrm{Hotab}_A \hand{\longrightarrow} \mathrm{Hotab}_{\hand{\Delta}} \]
is induced by $j^{*}_{ab}$ . Thus we get the wished for factorization of (6)
via $\mathrm{Hotab}_A$ , \uline{provided we can find an aspheric functor} (10). It
should not be hard moreover \inserted[typed]{to see} that the factorizing functor obtained
from (13), namely

(14) \qquad $\mathrm{Hotab}_A \hand{\longrightarrow} \mathrm{Hotab}$ \quad ,

does not depend up to canonical isomorphism on the choice of j , at
least in the case when A is a contractor, using the end remarks of
section 82 (p.272) concerning products of \xout{test}\inserted[typed]{aspheric} functors.

The question remains whether we can define (14) for (more or less)
any small category A, without \over{g}{h}aving to rely upon the \xout{\illegible{3}} existence of

\margin{6}

% Visual collation: both supplied image halves inspected; see CHECK_DECISIONS_134_7_8.tsv.
